% Two-page summary of technical_note_ocr_en.tex
% Build:  pdflatex summerize.tex
\documentclass[10pt,a4paper]{article}

\usepackage[utf8]{inputenc}
\usepackage[T1]{fontenc}
\usepackage{lmodern}
\usepackage[english]{babel}
\usepackage[a4paper,top=1.8cm,bottom=1.8cm,left=2cm,right=2cm]{geometry}
\usepackage{graphicx}
\usepackage{array}
\usepackage{enumitem}
\usepackage{float}
\usepackage{titlesec}

\newcolumntype{L}[1]{>{\raggedright\arraybackslash}p{#1}}

% booktabs-style rules, defined locally (booktabs is not in every MiKTeX setup)
\newcommand{\toprule}{\noalign{\hrule height 0.9pt}\noalign{\vskip 2.5pt}}
\newcommand{\midrule}{\noalign{\vskip 2.5pt}\noalign{\hrule height 0.4pt}\noalign{\vskip 2.5pt}}
\newcommand{\bottomrule}{\noalign{\vskip 2.5pt}\noalign{\hrule height 0.9pt}}

\titlespacing*{\section}{0pt}{7pt}{3pt}
\setlist[itemize]{itemsep=1pt,topsep=2pt,parsep=0pt,leftmargin=*}
\setlength{\parskip}{2pt}
\pagestyle{empty}

\begin{document}

% ------------------------------------------------------------------ header
\noindent
\begin{minipage}[c]{0.12\textwidth}
\includegraphics[width=1.6cm]{../ey.png}
\end{minipage}%
\begin{minipage}[c]{0.88\textwidth}
{\large\bfseries Intelligent OCR for Scanned Documents --- French / English / Arabic}\\[2pt]
{\small Summary of the technical note: rationale for the OCR engine selection --- July 2026}
\end{minipage}

\vspace{4pt}
\noindent\rule{\textwidth}{0.9pt}

% ------------------------------------------------------------------ 1
\section*{1. Problem}

Scanned administrative forms, reports and handwritten material must be
transcribed in \textbf{French, English and Arabic}, including mixed and
right-to-left content. Classical OCR (Tesseract) detects \emph{character
shapes} with no understanding of context, layout or language, and degrades
sharply on noisy scans, mixed scripts, tables and handwriting --- precisely the
cases at hand. A \textbf{Vision-Language Model} reads the image and produces
text within a single network, and that joint understanding is what makes
difficult documents tractable.

Three constraints shaped every decision: the reference test document is a
\textbf{strictly confidential} safeguarding form containing names and medical
details, so no document may leave the machine; the target hardware is a laptop
with an \textbf{NVIDIA RTX 4050 (6\,GB VRAM)}; and the solution must incur
\textbf{no recurring cost}.

% ------------------------------------------------------------------ 2
\section*{2. Solution}

A \textbf{local-first} chain: ingestion (PDFs rasterised locally, one image per
page) $\rightarrow$ preprocessing $\rightarrow$ VLM inference $\rightarrow$
Markdown structuring and language detection $\rightarrow$ optional quality
control $\rightarrow$ delivery. Preprocessing runs on \textbf{OpenCV}
(grayscale, non-local-means denoising, deskew, CLAHE contrast); binarisation is
applied only on the Tesseract path, since the vision model reads
contrast-enhanced grayscale better. PDFs are rasterised with \textbf{PyMuPDF}
at 300\,DPI --- the resolution at which fine print stays legible --- and read
page by page.

\begin{table}[H]
\centering\small
\begin{tabular}{@{}L{2.9cm}L{5.4cm}L{7.6cm}@{}}
\toprule
\textbf{Role} & \textbf{Model} & \textbf{Where it runs} \\
\midrule
Default engine & Qwen3-VL 4B instruct & Local, via Ollama, resident on the GPU \\
Second opinion (QA) & Qwen3-VL 8B instruct & Local, via Ollama \\
Classical baseline & Tesseract (\texttt{fra+eng+ara}) & Local, optional --- benchmarking and hybrid mode \\
Optional cloud & Llama 4 Scout (vision) & Groq API \\
\bottomrule
\end{tabular}
\end{table}

All four sit behind one abstraction, so engines can be swapped without touching
the processing chain. Three recognition modes run over it: \textbf{VLM vision}
(default), \textbf{hybrid} (Tesseract + VLM correction) and \textbf{Tesseract
only} (baseline).

% ------------------------------------------------------------------ 3
\section*{3. Key technical decisions}

\begin{itemize}
  \item \textbf{A VLM rather than classical OCR.} Classical OCR recognises
        glyphs; a VLM understands the page --- context, layout and language.
        Tesseract is kept as an optional baseline so the claim is measured, not
        asserted.
  \item \textbf{Qwen3-VL.} Strong multilingual performance including Arabic,
        built for document tasks (dense text, tables, structured forms,
        Markdown output), open weights, no licence fee.
  \item \textbf{4B rather than 8B --- measured, not assumed.} On the reference
        handwritten form the 4B scored \textbf{10/10 in 59\,s} against
        \textbf{9/10 in 135\,s} for the 8B: at 3.3\,GB it stays fully resident
        in 6\,GB of VRAM, while the 8B spills to the CPU. The smaller model is
        both faster and more complete on this hardware.
  \item \textbf{The \texttt{-instruct} tag is mandatory.} The base
        \texttt{qwen3-vl} tags run in \emph{thinking} mode: they spend the whole
        token budget on internal reasoning and return an empty transcription.
  \item \textbf{Context window at 8192 tripled the speed.} Each page is a
        separate call, so the window only needs to hold one page. Raising it
        further wastes VRAM on the KV cache and forces model layers onto the
        CPU --- about three times slower, for no accuracy gain.
  \item \textbf{Local (Ollama) as the default.} No document leaves the machine,
        no per-page fee, no rate limit, no internet needed --- and fast enough
        on real hardware.
\end{itemize}

% ------------------------------------------------------------------ 4
\section*{4. Quality control}

On a dense form the dominant failure mode is \textbf{omission}, not
hallucination: faint print, marginal notes, table cells and checkboxes get
silently dropped, and simply reading the page twice does not help --- the model
skips the same things twice. Three independent tools address this from
different angles, alongside CER/WER measurement where a ground truth exists.

\begin{table}[H]
\centering\small
\begin{tabular}{@{}L{3.4cm}L{9.2cm}L{3.1cm}@{}}
\toprule
\textbf{Tool} & \textbf{Principle} & \textbf{Cost} \\
\midrule
Recall verification &
  A second pass receives the same image \emph{plus its own draft}, under a
  prompt whose only job is to find what the draft missed. Handing the model its
  own output as a checklist turns open-ended transcription into a targeted
  search. & 2$\times$ model calls \\
Two-model comparison &
  The page is read by two independent models and every disagreement is flagged.
  The 4B and 8B have \emph{different} blind spots, so the diff surfaces real
  omissions a single pass would hide. Fully offline. & 2.5--3$\times$ \\
Format validation &
  Postcodes, phone numbers and emails are format-constrained, so a share of
  errors is catchable \textbf{with no ground truth at all}: structural,
  consistency and length checks, with an unambiguous repair suggested when a
  single OCR confusion substitution fixes the token. & negligible \\
\bottomrule
\end{tabular}
\end{table}

Format validation checks \textbf{form, not truth}: a well-formed but wrong value
passes clean. A clean report means ``no detectable format errors'', never ``no
errors''.

% ------------------------------------------------------------------ 5
\section*{5. Delivery and performance}

The OCR core is interface-agnostic. The primary deliverable is a
\textbf{full-stack web platform} --- React (Vite) single-page application,
FastAPI backend with JWT authentication and PBKDF2 password hashing, MongoDB
for a permanent searchable per-user document history --- backed by a Streamlit
prototype and a command-line interface that additionally exposes the three
quality tools above.

\begin{table}[H]
\centering\small
\begin{tabular}{@{}L{5.6cm}L{4.4cm}L{5.6cm}@{}}
\toprule
\textbf{Workload} & \textbf{Time} & \textbf{Notes} \\
\midrule
Printed page & 25--29\,s & RTX 4050, 300\,DPI, \texttt{num\_ctx} = 8192 \\
Dense handwritten page & 32--47\,s & \\
Three-page printed PDF & 86\,s & measured \\
Ten-page form & about 5 minutes & \\
First page after idle & +10\,s, one-time & model loads into VRAM \\
\bottomrule
\end{tabular}
\end{table}

Rule of thumb: \textbf{time $\approx$ pages $\times$ 30\,s}. Two parameters
trade speed against completeness --- the context window and the PDF
rasterisation resolution --- and both defaults deliberately favour
completeness.

% ------------------------------------------------------------------ 6
\section*{6. Known limitations}

\begin{itemize}
  \item Performance is tied to the presence of a GPU and its VRAM; without one,
        processing times become prohibitive.
  \item The figures come from one reference document on one machine: they
        indicate a trend, not a universal guarantee.
  \item Fidelity is not guaranteed at 100\%. Given personal and medical data,
        \textbf{human proofreading remains necessary} --- none of the three
        quality tools replaces it.
  \item Arabic handwriting remains the hardest case and accounts for most of the
        residual errors.
  \item The cloud path depends on a provider's model catalogue, which changes
        over time; the model identifier is therefore configurable.
  \item The format-validation rules are UK-specific and would need extending for
        other locales.
\end{itemize}

% ------------------------------------------------------------------ 7
\section*{7. Conclusion}

The system uses \textbf{Qwen3-VL-4B-Instruct} served locally via
\textbf{Ollama} as its default OCR engine, chosen for its multilingual
performance (notably Arabic), its handling of documents, tables and
handwriting, and --- critically --- the fact that it runs \textbf{fully offline
on a 6\,GB laptop GPU}. The 4B size was selected over the 8B after direct
measurement. A Groq / Llama 4 Scout cloud engine is retained behind the same
abstraction as a higher-capacity alternative and as a benchmark. Around that
engine the project delivers a complete chain --- preprocessing, page-by-page
PDF handling, structured Markdown output, three quality-control tools and a
multi-user web platform --- built to make human review short, not unnecessary.

\end{document}
